\documentclass[../../main.tex]{subfiles}
\begin{document}

空間内の点の集合 $\{ (x,y,z) \mid 0\leqq y,0\leqq z \}$ に含まれ，
原点 $O$ において $x$ 軸に接し，
$xy$ 平面と $45^\circ$ の傾きをなす，
半径1の円板 $C$ がある．
座標が $(0,0,2\sqrt2)$ の位置にある点光源 $P$ により，
$xy$ 平面上に投ぜられた円板 $C$ の影を $S$ とする．

\begin{figure}[H]
  \centering
  \tdplotsetmaincoords{70}{55}
  \begin{tikzpicture}[scale=1.4, >=stealth, every node/.style={font=\small}]
    \begin{scope}[tdplot_main_coords]
      % 座標軸
      \draw[->, thick, gray!80!black] (-1.5,0,0) -- (1.5,0,0) node[below left=1pt] {$x$};
      \draw[->, thick, gray!80!black] (0,-0.5,0) -- (0,3.2,0) node[right] {$y$};
      \draw[->, thick, gray!80!black] (0,0,-0.3) -- (0,0,3.2) node[above] {$z$};
      \node[below right=3pt] at (0,0,0) {$O$};

      % 円板 C（平面 y=z 内，中心 (0,√2/2,√2/2)，半径1，原点 O で x 軸に接する）
      \fill[pattern=north east lines, pattern color=blue!65, opacity=0.7]
      plot[domain=0:360, samples=60, variable=\t]
      ({cos(\t)}, {0.7071+0.7071*sin(\t)}, {0.7071+0.7071*sin(\t)});
      \draw[very thick, blue!85!black]
      plot[domain=0:360, samples=80, variable=\t]
      ({cos(\t)}, {0.7071+0.7071*sin(\t)}, {0.7071+0.7071*sin(\t)});
      \node[blue!90!black, above=2pt] at (0, 1.4142, 1.4142) {$C$};

      % 点光源 P
      \coordinate (P) at (0, 0, 2.8284);
      \node[circle, fill=red!85!black, inner sep=1.4pt] at (P) {};
      \node[above left=1pt, red!85!black] at (P) {$P(0,0,2\sqrt{2})$};

      % 光線
      \draw[dashed, red!70!black] (P) -- (0, 0, 0);
      \draw[dashed, red!70!black] (P) -- (0, 1.4142, 1.4142);
    \end{scope}
  \end{tikzpicture}
\end{figure}

\begin{enumerate}
  \item $S$ の輪郭を表す $xy$ 平面上の曲線の方程式を求めよ．
  \item 円板 $C$ と影 $S$ の間にはさまれ，光の届かない部分のつくる立体の体積を求めよ．
\end{enumerate}

\end{document}
